\section{Jatkuvat esitykset ja neuroverkkopohjainen kielimallinnus}\label{sec:mlp}

Siirtyminen diskreeteistä sanajonoista jatkuviin vektoriavaruuksiin ja neuroverkkopohjaiseen mallinnukseen on edistänyt merkittävästi kielimallien yleistyskykyä, sillä se mahdollistaa semanttisten suhteiden mallintamisen pelkän pintatason sanavastaavuuden sijaan \cite{bengio2003neural, mikolov2013efficient}.

Tämä luo teoreettisen pohjan tutkielman keskeiselle tutkimuskysymykselle, jossa tarkastellaan mallin kykyä tulkita luonnollisen kielen tiedontarpeita luotettavasti ja johdonmukaisesti merkityksen säilyttävistä kielellisistä variaatioista huolimatta.

\subsection{Upotusvektorit}

One-hot-esitys tarjoaa yksikäsitteisen tavan kuvata tokeneita, mutta sen
heikkoutena on, ettei se sisällä eksplisiittistä tietoa tokenien välisistä
suhteista. Jokaisen tokenin vektorikuvaus on ortogonaalinen kaikkiin muihin
nähden, jolloin esimerkiksi sanan \emph{kissa} suhde sanaan \emph{koira} on
rakenteellisesti sama kuin sen suhde sanaan \emph{tietokone}. Tämän rajoitteen
ylittämiseksi käytetään upotusvektoreita (engl.\ \emph{embedding vectors}),
joissa kukin token kuvataan tiheällä reaaliarvoisella vektorilla.

Olkoon edelleen \(\mathcal{V} = \{v_1, v_2, \dots, v_{|\mathcal{V}|}\}\) sanasto. Tällöin tokenille
\(v_j\) määritellään upotus
\[
  e(v_j) \in \mathbb{R}^{1 \times d},
\]
missä \(d \ll |\mathcal{V}|\) on upotustilan dimensio.

Upotusvektorit voidaan esittää myös matriisimuodossa. Olkoon
$$
E \in \mathbb{R}^{|\mathcal{V}| \times d}
$$
upotusmatriisi, jonka rivit vastaavat sanaston tokeneita. Kun
\(x(v_j) \in \mathbb{R}^{1 \times |\mathcal{V}|}\) on tokenin \(v_j\)
one-hot-esitys, sen upotus saadaan muodossa
$$
e(v_j) = x(v_j)E \in \mathbb{R}^{1 \times d}.
$$
Koska one-hot-vektori sisältää täsmälleen yhden nollasta poikkeavan komponentin, tämä kertolasku palauttaa suoraan tokenia \(v_j\) vastaavan upotusmatriisin rivin.

Upotusvektorien keskeinen etu liittyy siihen, että samankaltaisten tokenien
upotukset voivat sijaita lähellä toisiaan upotustilassa. Tällöin tokenien
geometrinen läheisyys voi heijastaa niiden semanttista tai syntaktista
samankaltaisuutta. Esimerkiksi sanat ``kuningas'' ja ``kuningatar'' voivat saada
upotukset, jotka ovat lähellä toisiaan, koska ne esiintyvät samankaltaisissa
konteksteissa ja jakavat merkityksellisiä piirteitä. Upotusvektorit tarjoavat
siten rikkaamman ja joustavamman esityksen tokenien välisistä suhteista kuin
one-hot-esitys, jossa kaikki tokenit ovat yhtä kaukana toisistaan.

Käytännön kielimallinnuksessa upotukset opitaan tavallisesti tehtäväkohtaisesti
osana laajempaa optimointia. Jos korpus on
\(\mathcal{W} = (w_1, w_2, \dots, w_{|\mathcal{W}|})\), voidaan jokainen tokeni \(w_t \in \mathcal{V}\)
korvata sitä vastaavalla upotusvektorilla \(e(w_t)\). Tällöin koko korpusta
voidaan tarkastella tiheiden vektorien jonona
$$
E_{\mathcal{W}} = (e(w_1), e(w_2), \dots, e(w_{|\mathcal{W}|})).
$$
Tämä esitystapa on erityisen hyödyllinen neuroverkoissa, joissa upotukset
muodostavat syötteen myöhemmille kerroksille ja joissa opittavat parametrit
voivat hyödyntää tokenien välisiä yhteyksiä yhteisen esitystilan kautta.

Upotusvektorit ovat myös yhteensopivia erikoistokenien kanssa. Esimerkiksi
aloitusmerkki \(\langle s \rangle\), lopetusmerkki \(\langle /s \rangle\) ja
tuntematon token \(\langle unk \rangle\) voidaan sisällyttää laajennettuun
sanastoon jolloin niille voidaan oppia omat upotuksensa samalla tavalla kuin muillekin
tokeneille. Näin myös sekvenssien rakenteelliset elementit voidaan käsitellä
yhtenäisesti muun sanaston kanssa.

Upotusvektorit muodostavat siten luonnollisen jatkeen one-hot-esitykselle.
Siinä missä one-hot kuvaa tokenin pelkkänä identiteettinä, upotusvektori tarjoaa
tiiviin ja optimoitavan esityksen, joka soveltuu tehokkaasti
neuroverkkopohjaisiin kielimalleihin.

\subsection{Upotusvektorien oppiminen Word2Vec-menetelmällä}

Upotusvektorit eivät tavallisesti ole ennalta määrättyjä, vaan ne estimoidaan
aineistosta optimointimenetelmän avulla. Tässä suhteessa yksi keskeisimmistä ja
vaikutusvaltaisimmista menetelmistä on Word2Vec \cite{mikolov2013efficient}. Menetelmän keskeinen tieteellinen kontribuutio ei
rajoitu siihen, että sanoille opitaan tiheitä vektoriesityksiä, sillä vastaavia
ajatuksia oli esitetty neuroverkkopohjaisessa kielimallinnuksessa jo aiemmin.
Word2Vecin varsinainen merkitys liittyy ennen kaikkea siihen, että se osoitti
tällaisten representaatioiden oppimisen olevan mahdollista laskennallisesti
erittäin tehokkaalla tavalla myös hyvin suurissa korpuksissa. Lisäksi se toi
esiin, että kontekstuaaliseen ennustamiseen perustuva oppiminen voi johtaa
vektoriavaruuksiin, joissa ilmenee kielellisesti mielekkäitä semanttisia ja
syntaktisia säännönmukaisuuksia.

Word2Vec perustuu distributionaaliseen hypoteesiin, jonka mukaan sanan merkitys
voidaan päätellä sen käyttöympäristöjen perusteella. Toisin sanoen sanat, jotka
esiintyvät samankaltaisissa konteksteissa, saavat oppimisprosessissa
samankaltaiset vektoriesitykset. Tämä periaate on laajasti koko modernin
kieliteknologian taustalla, mutta Word2Vec teki siitä käytännöllisen ja
skaalautuvan laskennallisen menetelmän. Menetelmä osoitti, että diskreettien
symbolien semanttista rakennetta voidaan mallintaa jatkuvassa vektoriavaruudessa
pelkästään paikallisiin konteksteihin perustuvan ennustetehtävän avulla ilman
käsin suunniteltuja kielellisiä piirteitä.

Olkoon korpus \(\mathcal{W} = (w_1, w_2, \dots, w_{|\mathcal{W}|})\), missä \(w_t \in \mathcal{V}\) on
sanaston \(\mathcal{V}\) sana ajanhetkellä \(t\). Word2Vec-malleissa tarkastellaan
tyypillisesti kunkin sanan ympärille määriteltyä paikallista kontekstia, jonka
laajuus määräytyy ikkunaparametrilla \(c \in \mathbb{N}\). Tällöin ajanhetken
\(t\) kontekstiksi voidaan määritellä
$$
C_t = \{w_{t-c}, \dots, w_{t-1}, w_{t+1}, \dots, w_{t+c}\},
$$
olettaen, että indeksit pysyvät korpuksen rajojen sisällä. Mikolov ym.\
\cite{mikolov2013efficient} tarkastelevat kahta vaihtoehtoista
perusarkkitehtuuria: jatkuvien sanojen pussimallia
(\emph{continuous bag-of-words}, CBOW) ja Skip-gram-mallia.

CBOW-mallissa pyritään ennustamaan tarkasteltava keskussana sen
kontekstisanojen perusteella. Tavoitteena on siten approksimoida
todennäköisyyttä
$$
p(w_t \mid w_{t-c}, \dots, w_{t-1}, w_{t+1}, \dots, w_{t+c}).
$$
Kontekstin sanat projisoidaan ensin vektoriesityksiksi, minkä jälkeen nämä
esitykset yhdistetään tavallisesti summaamalla tai keskiarvoistamalla.
Muodostettua kontekstivektoria käytetään tämän jälkeen kohdesanan
ennustamiseen koko sanaston yli. CBOW-mallin etuna on sen laskennallinen
tehokkuus erityisesti suurilla aineistoilla, koska useiden kontekstisanojen
informaatio voidaan tiivistää yhdeksi esitykseksi ennen ennustevaihetta
\cite{mikolov2013efficient}.

Skip-gram-mallissa ennustesuunta on päinvastainen. Siinä tavoitteena on
ennustaa keskussanan avulla sitä ympäröivät kontekstisanat. Tällöin
maksimoidaan todennäköisyydet
$$
p(w_{t+j} \mid w_t),
\qquad
-c \le j \le c,\; j \ne 0.
$$
Koko korpuksen tasolla vastaava tavoitefunktio voidaan kirjoittaa muodossa
$$
\sum_{t=1}^{|\mathcal{W}|} \sum_{\substack{-c \le j \le c \\ j \ne 0}}
\log p(w_{t+j} \mid w_t).
$$
Skip-gram-malli osoittautuu erityisen käyttökelpoiseksi silloin, kun tavoitteena
on oppia informatiivisia esityksiä myös harvinaisemmille sanoille, sillä kukin
havaittu sana toimii useiden kontekstisanojen ennustamisen lähtökohtana
\cite{mikolov2013efficient}.

Olkoon \(w_I\) syötesana ja \(w_O\) kohdesana. Jos syötesanaa vastaava
rivivektoriesitys on \(v_{w_I} \in \mathbb{R}^{1 \times d}\) ja kohdesanaa
vastaava ulostulorivivektori on \(u_{w_O} \in \mathbb{R}^{1 \times d}\),
voidaan todennäköisyys \(p(w_O \mid w_I)\) esittää softmax-funktion avulla
muodossa
$$
p(w_O \mid w_I)
=
\frac{\exp(v_{w_I} u_{w_O}^{\top})}
{\sum_{w \in \mathcal{V}} \exp(v_{w_I} u_w^{\top})}.
$$

Tässä nimittäjän summa ulottuu koko sanastoon, mikä tekee täsmällisestä
optimoinnista laskennallisesti raskasta suurilla sanastoilla. Juuri tämän
ongelman ratkaiseminen on yksi Word2Vecin keskeisistä ansioista. Menetelmä teki
näkyväksi sen, että suuren sanaston ulostuloprojektio muodostaa
vektoriesitysten oppimisessa keskeisen laskennallisen pullonkaulan, ja esitti
tähän käytännöllisiä approksimaatioita. Näistä erityisen merkittävä on
negatiivinen otanta (\emph{negative sampling}), jossa koko sanaston yli
määritellyn softmax-todennäköisyyden optimoinnin sijasta ratkaistaan joukko
binäärisiä luokittelutehtäviä.

Olkoon \((w_I, w_O)\) havaittu syöte--kohdesana-pari, jossa \(w_O\) esiintyy
sanan \(w_I\) kontekstissa. Negatiivisessa otannassa valitaan lisäksi
\(K\) negatiivista sanaa \(w_1^{-}, \dots, w_K^{-}\) jakaumasta
\(P_n(w)\), joka on tyypillisesti jokin sanaston sanoille määritelty
näytteenottojakauma. Tällöin mallin tavoitteena ei ole enää optimoida
todennäköisyyttä \(p(w_O \mid w_I)\) koko sanaston yli, vaan maksimoida
havaittua paria vastaava binäärinen logistinen tavoite
$$
\log \sigma\!\left(v_{w_I} u_{w_O}^{\top}\right)
+
\sum_{k=1}^{K}
\log \sigma\!\left(-v_{w_I} u_{w_k^{-}}^{\top}\right),
$$
missä \(\sigma(x) = \frac{1}{1 + e^{-x}}\) on logistinen sigmoidifunktio.
Ensimmäinen termi kasvattaa havaittuun positiiviseen sana--konteksti-pariin
liittyvän pistetulon \(v_{w_I} u_{w_O}^{\top}\) arvoa, kun taas jälkimmäinen
termi pienentää satunnaisesti valittujen negatiivisten sanojen
\(v_{w_I} u_{w_k^{-}}^{\top}\) arvoja. Näin malli oppii erottamaan aidot
kontekstisanat satunnaisista sanoista ilman, että jokaisella oppimisaskeleella
tarvitsisi normalisoida jakaumaa koko sanaston yli.

Koko aineiston yli tavoite voidaan esittää odotusarvomuodossa
$$
\sum_{(w_I, w_O) \in D}
\left[
  \log \sigma\!\left(v_{w_I} u_{w_O}^{\top}\right)
  +
  \sum_{k=1}^{K}
  \mathbb{E}_{w_k^{-} \sim P_n}
  \log \sigma\!\left(-v_{w_I} u_{w_k^{-}}^{\top}\right)
\right],
$$
missä \(D\) on havaittujen syöte--konteksti-parien joukko ja \(K\) on negatiivisten otantojen määrä. Tällä tavoin
alkuperäinen moniluokkainen softmax-optimointi approksimoidaan joukolla
logistisia regressiotehtäviä, joissa erotellaan positiiviset havaintoparit
negatiivisesti näytteistetyistä pareista. Laskennallinen etu syntyy siitä, että
kunkin havaintoparin kohdalla päivitetään vain positiivista kohdesanaa ja
\(K\):ta negatiivista sanaa vastaavat parametrit koko sanaston kattavan
päivityksen sijasta.

Word2Vec-mallissa opitaan kaksi parametrimatriisia: syötevektorit ja
ulostulovektorit. Jokaiselle sanaston sanalle \(v_j \in \mathcal{V}\) assosioituu siten
ainakin yksi tiheä vektoriesitys. Oppimisen jälkeen sanan lopullisena
esityksenä käytetään tavallisesti syötevektoria tai vaihtoehtoisesti syöte- ja
ulostulovektoreiden yhdistelmää. Näin saadut upotukset heijastavat korpuksessa
havaittuja yhteisesiintymisrakenteita siten, että samankaltaisissa
käyttöympäristöissä esiintyvät sanat sijoittuvat tyypillisesti lähelle toisiaan
upotustilassa.

Mikolov ym.\ \cite{mikolov2013efficient} osoittavat lisäksi, että Word2Vecin
oppimat vektoriesitykset eivät ainoastaan ryhmittele semanttisesti läheisiä
sanoja, vaan niihin syntyy myös lineaarisia relaatiota vastaavia
säännönmukaisuuksia. Sanavektorien erotukset voivat vastata esimerkiksi
sukupuoleen, lukuun tai taivutusmuotoon liittyviä semanttisia ja syntaktisia
eroja, mikä ilmenee analogiatehtävissä. Tämä havainto on menetelmän kannalta
erityisen merkittävä, sillä se viittaa siihen, että kontekstuaaliseen
ennustamiseen perustuva optimointi tuottaa emergentisti sellaisia
vektoriavaruuden rakenteita, jotka eivät ole eksplisiittisesti koodattuja
malliin, mutta jotka vastaavat kielellisesti tulkittavia suhteita. Näin
Word2Vec osoitti, että jakaumallinen oppiminen ei tuota ainoastaan
samankaltaisuusmittaa, vaan voi synnyttää myös abstraktimpia relaatioita
mallintavia geometrisia rakenteita.

Vaikka Word2Veciä ei nykyisissä suurissa kielimalleissa tyypillisesti käytetä
sellaisenaan, sen vaikutus näkyy edelleen usealla
tasolla. Ensinnäkin se vakiinnutti distributionaalisen hypoteesin
neuroverkkopohjaisen kielimallinnuksen keskeiseksi oppimisperiaatteeksi.
Toiseksi se toi esiin suuren sanaston ulostuloprojektion laskennallisen
haasteen sekä popularisoi tehokkaita approksimointimenetelmiä, kuten
negatiivisen otannan. Kolmanneksi se osoitti vakuuttavasti, että
yksinkertainenkin kontekstuaalinen ennustetehtävä voi johtaa
vektoriavaruuksiin, joissa semanttiset ja syntaktiset suhteet ilmenevät
lineaarisina rakenteina. Näistä syistä Word2Vecia voidaan pitää tärkeänä
siirtymävaiheena klassisten jakaumallisten sanamallien ja myöhempien
syvien neuroverkkojen, mukaan lukien nykyaikaiset suuret kielimallit,
välillä.
